\scanpage{101}
% archive sheet number: 103

2) \uline{$E'_a$} Donc $\pi$ engendré par $u,v$ avec $\sigma u=u,\ \sigma v=u-v$

\[
\begin{aligned}
&\sigma(x,y)=\sigma(xu+yv)=xu+y(u-v)=(x+y)u-yv.\\
&\sigma(x,y)-(x,y)\unsure{\equiv|=}yu-2yv=y(u-2v),
\end{aligned}
\]

donc la condition est $y\in\mathbb Z$ ; donc on trouve \deleted{que} les points de $\mathbb R\unsure{u|e_0|u_+}$ mod $\pi$, donc dans ce cas $X(\mathbb R)$ est connexe.

\margin{$9.2$}

\uline{NB} Pour bien faire, il faudrait aussi classifier les courbes elliptiques réelles inhomogènes, i.e. les torseurs (sur $\operatorname{Spec}\mathbb R$) sous les \inserted{schémas en} groupes elliptiques $E_a,E'_a$. Les torseurs non triviaux fournissent les surfaces conformes \deleted{sans} bord non orientables de genre $1$. La suite exacte de $\sigma$-modules (où $\gamma=\operatorname{Gal}(\mathbb C/\mathbb R)\simeq\{\pm1\}$)

\diagram{\begin{tikzcd}[ampersand replacement=\&,column sep=small]
0\arrow[r,dash]\&\pi\arrow[r,dash]\&V\arrow[r,dash]\&E(\unsure{\mathbb C|\deleted{\mathbb C}})\arrow[r,dash]\&0
\end{tikzcd}}

fournit

\[
H^1(\gamma,E(\mathbb C))\simeq H^2(\gamma,\pi)\simeq\pi^\gamma/(1+\sigma)\pi=\pi_+/(1+\sigma)\pi
\]

1) Cas de $E_a$. On a $\pi\simeq\pi_+\oplus\pi_-$, d'où $\pi^\gamma=\pi_+$

\[
H^1(\mathbb R,E_{\mathbb R})\simeq H^1(\gamma,E(\mathbb C))\simeq H^2(\gamma,\pi)\simeq
\underset{\Vert\,\mathbb Z/2\mathbb Z}{\underline{H^2(\gamma,\pi_+)}}\oplus
\underset{\Vert\,0}{\underline{H^2(\gamma,\pi_-)}}
\]

Il y a exactement un élément non nul dans $H^1$, qui est d'ordre $2$

2) Cas de $E'_a$. On a $H^1(\mathbb R,E_{\mathbb R})\simeq0$.

Donc les seules courbes elliptiques vraiment inhomogènes sur $\mathbb R$ sont les $E_a^{!}=$ \inserted{les} torseurs \deleted{\unsure{tr|ts}} non triviaux (d'ordre $2$) sous $E_a$, dépendant du paramètre $a\in\mathbb R^{*+}$.
